\chapter{Machine Consciousness}

\section{Introduction}

Consciousness has remained one of the most contested and enigmatic concepts within both philosophy and modern science. Despite centuries of inquiry, clear definitions and consensus on its origins, mechanisms, and significance are still elusive \cite{francken2022consciousness,reggia2013rise}. This ambiguity extends not only to the foundational question of what consciousness is, but also to debates about whether it is exclusively biological or, as some propose, could potentially arise in artificial systems.

Broadly, consciousness is often reduced to brain activity in scientific discourse. However, leading voices---both within and outside neuroscience---argue this reduction overlooks critical dimensions of conscious experience, particularly its subjective and creative qualities \cite{garridomerchan2024machine}. The brain may enable sensory perceptions and bodily awareness, yet it does not suffice to claim consciousness is simply a product of neural computation. Such reductionist assumptions are now being challenged by physicists and philosophers positing that consciousness is fundamentally irreducible and cannot be reconstructed by algorithmic or computational means \cite{faggin2024irreducible,penrose1990emperors}. These arguments form the cornerstone of a critical stance toward claims of Artificial Consciousness (AC).

Nevertheless, rapid progress in artificial intelligence and information processing has revitalized philosophical and empirical explorations into the possibility of Machine Consciousness (MC). These two terms, AC and MC, might be confused, but AC refers to a phenomenon, whereas MC concerns the application of AC to a particular kind of entity, namely a machine. First and foremost, a machine here refers to an information-processing system or computing device. Researchers in neuroscience, philosophy of mind, and computer science have sought to clarify both what is meant by \emph{consciousness} and under what (if any) circumstances an artificial entity might possess it \cite{butlin2023consciousness}. Despite optimistic claims, mainstream scientific consensus acknowledges that current AI systems display, at best, advanced forms of self-monitoring or adaptive behavior---what might be called \emph{access consciousness} or \emph{self-awareness}---while the presence of true subjective, phenomenal consciousness (``what it is like to be a system,'' \cite{nagel1974like}) remains speculative and unsupported by empirical evidence \cite{francken2022consciousness,garridomerchan2024machine,reggia2013rise}.

This chapter begins by critically surveying the philosophical foundations and neuroscientific theories of consciousness. It discusses major conceptual currents in the field, including both computational and non-computational (epiphenomenal, quantum) theories. Next, the relationship between self-awareness, adaptive behavior, and consciousness is explored in living and artificial systems \cite{graziano2022framework}. Special attention is paid to the ongoing debate about whether information processing alone can serve as the basis for conscious experience, or whether the irreducibility of consciousness poses a fundamental limit to the reach of artificial cognition \cite{gamez2008progress,searle1980minds}.

Finally, the chapter adopts a critical and epistemologically rigorous approach. While acknowledging that computer systems can, in principle, simulate a wide range of intelligent and adaptive behaviors, it will argue, in alignment with Faggin's irreducibility thesis and recent philosophical critiques, that algorithmic architectures---regardless of complexity, quantum extension, or learning---cannot replicate creative, subjective experience \cite{faggin2024irreducible,garridomerchan2024machine}. The implications of this stance reach into the ongoing discourse surrounding artificial general intelligence (AGI), autonomous moral agents, and the prospects of reaching technological singularity.

Studying machine consciousness remains a fertile ground for clarifying the ontology of consciousness, refining theories in philosophy of information, and illuminating what may ultimately distinguish creative, conscious beings from all forms of artificial information processing. By examining these fundamental distinctions, science and philosophy can together approach, but perhaps never fully resolve, the riddle of conscious machines.

\begin{figure}[htbp]
\centering
\includegraphics[width=0.99\textwidth]{media/image1.png}
\caption{Concept map of philosophical theories of consciousness. In the first layer are ontological concepts divided into reductionist (darker bottom) and non-reductionist (lighter top). The second layer consists of theories, models, and critical arguments implied by these stances. The third layer groups critiques of computationalism and formal approaches which leads to the denial of functionalism and rejection of MC hypothesis. The fourth layer indicates the implication of the theories for MC: which views reject, accept, or leave unresolved the possibility of conscious machines.}
\label{fig:consciousness_map}
\end{figure}

\section{Philosophical Foundations of Consciousness}

This section reviews essential philosophical debates and definitions surrounding consciousness. The discussion begins with phenomenological perspectives (Nagel, Chalmers), addresses questions of qualia, intentionality, and \emph{the hard problem}, and distinguishes between \emph{phenomenal} and \emph{access consciousness}. Competing philosophical traditions (dualism, materialism, functionalism, panpsychism) and key theorists are reviewed. The concept map in Figure~\ref{fig:consciousness_map} is a summary of the topics included in this investigation and their implications for the \emph{hypothesis of conscious machines} .

\subsection{Defining Consciousness: Historical and Disciplinary Perspectives}

Consciousness has persisted as a central yet perplexing subject within philosophy from antiquity to the modern era, inspiring diverse schools of thought about its origins, properties, and role in reality \cite{francken2022consciousness,reggia2013rise}. Though its everyday use refers broadly to self-awareness or being awake, philosophical scrutiny splits consciousness into multiple, highly nuanced constructs---most centrally, phenomenal consciousness and access consciousness \cite{block1995confusion,rosenthal2009concepts}.

\begin{itemize}
\item \emph{Phenomenal consciousness} refers to subjective experience, or qualia---the intrinsic ``what it is like'' aspect of sensation, emotion, and thought \cite{levine1983materialism,nagel1974like}.
\item \emph{Access consciousness} involves information in the mind available for reasoning, reporting, or guiding action, irrespective of any accompanying subjective feeling \cite{block1995confusion}.
\end{itemize}

These distinctions profoundly shape both philosophical and scientific debates, as the former addresses the irreducibility and qualitative nature of mental life, while the latter grounds conscious processes within observable cognitive functioning and behavioral outputs \cite{garridomerchan2024machine,reggia2013rise}. Conceptual ambiguities remain a major obstacle: a recent survey of leading consciousness researchers revealed enduring disagreement about definitions, methodologies, and even what constitutes an explanation for consciousness \cite{francken2022consciousness}.

\subsection{The ``Hard Problem'' and Qualia}

David Chalmers \citeyear{chalmers1996conscious} famously divided consciousness research into \emph{easy} and \emph{hard problems}. Easy problems involve functional mechanisms and information processing---tasks like perception, attention, and verbal reporting---widely seen as scientifically tractable \cite{chalmers1995facing,dennett1991consciousness}. \emph{The hard problem}, however, confronts the emergence of subjective experience: why does processing sensory data or integrating information result in a feeling, and why does it feel like anything at all \cite{levine1983materialism,nagel1974like}? This \emph{explanatory gap} remains one of the most intractable questions in philosophy \cite{francken2022consciousness,levine1983materialism}.

Qualia---the felt qualities of red, pain, or joy---are deeply private, ineffable, and, many argue, not reducible to third-person scientific description \cite{nagel1974like,schwitzgebel2016phenomenal}. The \emph{zombie argument} \cite{chalmers1995facing,kirk1974zombies} and \emph{knowledge argument} \cite{jackson1982epiphenomenal} further illustrate this conceptual challenge: intelligent behavior and information processing do not necessitate experience, and having all physical knowledge about color vision does not guarantee an understanding of the experience of \emph{seeing red} \cite{francken2022consciousness}.

\subsection{Major Philosophical Theories}

\subsubsection{Dualism}

Dualism posits the existence of two distinct realms---physical and mental---making consciousness a non-material property or substance \cite{churchland1999engine,descartes1641meditations,eccles1994self}. Despite its intuitive appeal and prevalence in popular thought, dualism faces substantial objections in scientific circles---chiefly the challenge of explaining interaction between physical and non-physical substances and the absence of empirical support \cite{reggia2013rise}.

\subsubsection{Monism}

Monist theories---idealism and physicalism---insist reality is fundamentally unified, but disagree on its essential nature \cite{reggia2013rise}. \emph{Idealism} views consciousness as foundational, with the physical universe as a product of mind \cite{goswami1993self,hutto1998presence}. This stance faces methodological difficulties and risks veering into solipsism. \emph{Physicalism} treats consciousness as an emergent or reducible property of matter and energy \cite{churchland1999engine,searle2004turn}.

Materialist perspectives branch out: (1) \emph{Identity theory:} Mind and brain states are one and the same, each mental state corresponding to a physical/neuronal process. (2) \emph{Philosophical behaviorism:} Reduces mental states to observable behaviors, circumventing questions about subjective experience. (3) \emph{Eliminative materialism:} Predicts future neuroscience will replace folk psychological concepts with more precise scientific ones, dissolving \emph{consciousness} as currently conceived.

\subsubsection{Functionalism}

Functionalism argues mental states are defined by their causal roles, not their material substrate---analogous to a computer's logic being independent of its hardware \cite{churchland1999engine,searle2004arguments}. This allows for \emph{multiple realizability}: the possibility that consciousness could, in principle, emerge in artificial systems with appropriate functional organization. Functionalism assumes that conscious mental states can be instantiated in any system---biological or artificial---that realizes the same functional relationships \cite{reggia2013rise}. Yet, critics claim functionalism cannot fully resolve \emph{the hard problem} nor address why subjective experience accompanies these functions \cite{garridomerchan2024machine,levine1983materialism}.

\subsubsection{Emergentism and Panpsychism}

Emergentism asserts that consciousness arises as an emergent property from sufficiently complex physical systems, particularly brains organized in specific ways \cite{chalmers2000conscious,francken2022consciousness}. Emergentism occupies a middle ground between reductive physicalism and dualism by maintaining that consciousness is dependent on physical processes yet not straightforwardly reducible to them, instead arising when neural activity and information reach specific organizational thresholds \cite{chalmers2000conscious}. In this view, consciousness cannot be fully predicted or explained by studying individual neural components in isolation; it is thought to emerge from holistic interactions and global integration of information, often connected to large-scale network dynamics or integrative processing models \cite{francken2022consciousness,tononi2016integrated}. Emergentism thereby bridges physicalist and non-reductive approaches, acknowledging both a material basis and irreducible complexity. However, it remains vulnerable to the \emph{explanatory gap} critique: why does physical complexity, however sophisticated, necessarily produce subjective experience rather than merely more complex behavior \cite{chalmers1995facing,francken2022consciousness,levine1983materialism}?

Panpsychism is a philosophical position asserting that consciousness (or mind-like properties) is a fundamental and ubiquitous feature of the universe---present, in some form, in all matter or all physical systems, rather than being unique to biological brains or emerging only at a certain level of complexity \cite{chalmers2016combination,strawson2008realistic}. The view is typically contrasted with physicalist or functionalist theories that locate consciousness exclusively in biological organisms (especially humans and some animals) or treat it as arising only when systems achieve sufficient computational sophistication \cite{reggia2013rise}. As such, panpsychism represents a radical departure from conventional neuroscientific and philosophical approaches that treat consciousness as either emergent from complexity or as an epiphenomenon of specific neural processes \cite{chalmers2016combination,francken2022consciousness}.

In other words, panpsychism posits that consciousness is not an emergent property arising from specific arrangements of matter, but rather an intrinsic, fundamental property of reality, analogous to mass or charge as basic properties in physics. The theory comes in several variants. Cosmopsychism claims that consciousness is primarily a property of the universe as a whole, with local conscious subjects derived from a cosmic mind, whereas constitutive panpsychism holds that individual entities---atoms, particles, simple organisms---possess rudimentary experiential properties that, when appropriately combined, yield higher-order conscious experiences \cite{chalmers2016combination}. On such views, subjectivity is not reserved for humans or animals with complex brains; even fundamental particles or very simple systems might instantiate minimal, primitive forms of ``what-it-is-likeness,'' though these would be vastly unlike human consciousness \cite{strawson2008realistic}.

Panpsychism introduces a significant challenge and alternative perspective to machine-consciousness debates. If consciousness is truly fundamental to all matter, then sophisticated AI systems---composed of physical substrates such as silicon and electronic circuits---might, in principle, participate in consciousness to some degree, just as biological brains do \cite{butlin2023consciousness,reggia2013rise}. At the same time, panpsychism faces substantial philosophical difficulties. The combination or binding problem remains unresolved: how do many simple, putatively conscious constituents (particles, micro-systems) combine into unified, macro-level consciousness with a single coherent subject of experience? Attributing consciousness to particles or atoms also appears to strain parsimony and empirical plausibility, and current panpsychist frameworks struggle to specify when and how consciousness in simple systems emerges or combines into richer forms. Critics argue that panpsychism risks being unfalsifiable and leaves \emph{the hard problem} of consciousness fundamentally unsolved, merely relocating the mystery from \emph{why consciousness arises at all} to \emph{why combinations of simple conscious entities produce complex consciousness} \cite{chalmers2016combination,francken2022consciousness}. Moreover, by treating consciousness as ubiquitous, panpsychism can blur crucial distinctions between genuine phenomenal consciousness and mere information processing, potentially conflating structural or functional complexity with experiential phenomena in ways that undermine rigorous criteria for machine consciousness \cite{garridomerchan2024machine,reggia2013rise}.

Notably, both emergentist and panpsychist perspectives intersect with contemporary irreducibility arguments that emphasize the limits of purely computational or information-processing accounts of consciousness. Faggin \citeyear{faggin2024irreducible}, for example, defends a form of \emph{irreducibility} according to which consciousness is a fundamental, creative field-like aspect of reality that cannot be captured by algorithmic structures alone, thereby resonating with panpsychist and non-reductive intuitions while sharpening the challenge posed to strong computationalist claims about artificial consciousness \cite{faggin2024irreducible,penrose1990emperors}.

\subsection{Contemporary Philosophical Currents}

Philosophy of mind has moved toward integration with neuroscience and information theory, producing new hybrid paradigms but perpetuating foundational disagreements: (1) \emph{Predictive processing theory} posits a brain constantly running top-down models, explaining experience as an inference from prediction-error minimization \cite{seth2022theories}. (2) \emph{Global Workspace Theory} \cite{baars2001conscious}: Consciousness as global brain broadcast, allowing information integration and flexible control. (3) \emph{Higher-Order Theories:} Mental states become conscious when represented by a higher-order thought or model \cite{rosenthal2009concepts}. (4) \emph{Integrated Information Theory (IIT):} Proposes consciousness correlates with maximal integration of information across a system, mapped by mathematically defined measures \cite{tononi2016integrated}.

Despite technical advances and empirical sophistication, many philosophers and scientists believe a complete theory must transcend functional and computational approaches to address the intrinsic, subjective dimensions of experience \cite{chalmers2000conscious,garridomerchan2024machine}.

\subsection{Eastern and Non-Western Philosophical Perspectives}

Beyond Western analytic philosophy, Eastern traditions offer rich and nuanced understandings of consciousness that emphasize its pervasive, dynamic, and foundational nature. Indian Vedantic and Buddhist philosophies, for example, frame consciousness as a fundamental reality rather than an emergent or derived property \cite{vonbraun2021consciousness}. Buddhism, particularly Yogacara and Madhyamaka schools, interpret consciousness not as a static substance but as a continuous flow or process lacking inherent essence. These traditions often reject the subject-object dualism central to many Western accounts, instead focusing on direct experiential insight \cite{siderits2007buddhism}.

While differing ontologically and metaphysically from Western physicalist and dualist views, Eastern philosophies contribute valuable perspectives on the unity of consciousness, the role of experience, and the nature of selfhood \cite{vonbraun2021consciousness}. Their rich phenomenological accounts have inspired contemporary theories in cognitive science, particularly in attention and self-awareness studies \cite{graziano2022framework}.

\subsection{Influential Thought Experiments and Conceptual Tools}

Philosophy of consciousness has developed numerous thought experiments to illustrate and challenge prevailing theories. The ``philosophical zombie''---a being physically identical to a human yet devoid of subjective experience---questions whether physicalism can fully capture consciousness \cite{chalmers1996conscious,kirk1974zombies}. \emph{Mary's Room Argument} \cite{jackson1982epiphenomenal} challenges materialism by imagining a neuroscientist who knows all physical facts about color but has never experienced it. \emph{Chinese Room Argument} \cite{searle1980minds} critiques functionalism by arguing that syntactic symbol manipulation, however complex, does not entail understanding or consciousness. These experiments highlight aspects of consciousness---subjectivity, experience, understanding---that resist reduction to functional or physicalist explanations \cite{francken2022consciousness}.

\subsection{Epistemology and Ontology of Consciousness}

Consciousness raises profound epistemological questions: how can subjective experience be accessed, studied, and validated? Phenomenology emphasizes first-person methods and descriptions to complement third-person science \cite{gallagher2022phenomenology}. Epistemic challenges persist, as the subjective nature of experience resists objective measurement. Ontologically, whether consciousness is fundamental or emergent remains debated---with positions ranging from physicalist reductionism to panpsychism and dual-aspect monism \cite{chalmers2016combination,strawson2008realistic}.

\subsection{The Creative, Irreducible, and Non-computational Nature of Consciousness}

Recent philosophical currents have advanced the argument that consciousness is not merely a complex product of physical or computational processes, but rather an irreducible, creative, and potentially non-computable phenomenon \cite{faggin2024irreducible,lipton2016biology,penrose1990emperors}. Federigo Faggin, for example, draws from quantum mechanics and subjective agency to argue that consciousness might be fundamental---a field-like phenomenon governing information and physical reality, not emerging from them \cite{faggin2024irreducible}. Penrose \citeyear{penrose1990emperors} similarly contends that human consciousness harbors aspects of non-computability---rooted in quantum physics and evident in mathematical intuition---that no purely algorithmic process can replicate. These ideas echo Albert Einstein's view that the ``field is the sole governing agency of the particle,'' suggesting a holistic substrate for conscious experience that transcends conventional informational theories.

In alignment with Bruce Lipton's \citeyear{lipton2016biology} insights on epigenetics and creative agency, these views complicate the reductionist standpoint. If consciousness is indeed capable of shaping biological outcomes via environmental feedback and beliefs, this implies causal powers beyond mere physical or algorithmic determinism---a challenge for both classical and contemporary models of machine consciousness.

\subsection{Machine Consciousness: Philosophical Objections and Simulated Phenomenology}

The question arises: even if a machine can simulate the functions of conscious beings, including advanced self-awareness and adaptive information processing, can it ever partake in real subjective experience? Searle's \emph{Chinese Room} argument contends that syntactic manipulation of symbols, no matter how complex, will not create understanding or experience \cite{searle1980minds}. Garrido-Merchán \citeyear{garridomerchan2024machine} and others extend such arguments to assert that present-day claims about MC severely overreach: current AI may model, learn, and self-describe, but there is no evidence it hosts qualia or creativity.

Furthermore, Integrated Information Theory \cite{tononi2016integrated}, Global Workspace Theory \cite{baars2001conscious}, and related models provide metrics or operational tests for the substrate-level correlates of consciousness. However, critics argue that these frameworks may only describe necessary organizational features for consciousness, not sufficient mechanisms for subjective or creative experience \cite{francken2022consciousness}. Philosophers note that quantifying information integration or workspace architectures cannot by itself bridge \emph{the explanatory gap} or fully address \emph{the hard problem} \cite{levine1983materialism}.

This critique is amplified by the irreducibility and non-computability arguments outlined in quantum-informed and phenomenological accounts: if consciousness is an ontological primitive or grounded in fundamental ``fields'' or non-local processes, no information-processing architecture---classical or quantum---could suffice for its implementation \cite{faggin2024irreducible,garridomerchan2024machine,penrose1990emperors}.

\subsection{Towards a Synthesis: What Consciousness Demands}

In summary, the philosophical foundations of consciousness are marked by rigorous debate between:

\begin{itemize}
\item Reductive and non-reductive theories of mind,
\item Functionalist/algorithmic explanatory models and holistic/field-like ontologies,
\item Western analytic, phenomenological, and Eastern monist traditions,
\item Empirical, information-theoretic, and quantum-mechanical perspectives.
\end{itemize}

These~positions inform~the subsequent evaluation of artificial and machine consciousness, notably highlighting (1) the difficulty (or impossibility) in moving from simulation of conscious functions to real experience, (2) the essential creative, causal, and irreducible properties attributed to consciousness in both philosophical and biological models, and (3) the risk of category error when equating self-aware computation with conscious agency or qualia. By critically reviewing these foundational arguments, the stage is set for a detailed inquiry into how, or even if, computational systems could cross the threshold from simulating intelligence and self-modeling to hosting genuine consciousness or creativity. This synthesis is central to ongoing debates in the philosophy of mind, cognitive neuroscience, and the ethics of AI development \cite{butlin2023consciousness,francken2022consciousness}.

\section{Theories of Consciousness: Phenomena and Mechanisms}

This section analyzes the most prominent scientific and philosophical theories of consciousness, with attention to how they interpret the emergence of subjective experience and what this implies for both natural and artificial consciousness. Building on previous philosophical distinctions---particularly the challenges of subjective qualia and \emph{the explanatory gap}---we now focus on the central theoretical frameworks proposed to explain conscious phenomena, and their potential (or limitations) regarding machine implementation.

\subsection{Epiphenomenal and Quantum Approaches: Consciousness as Irreducible Phenomenon}

\emph{Epiphenomenal theories} posit that consciousness may be an emergent or accompanying phenomenon of physical processes, but not reducible to those processes nor causally efficacious in itself. Among the most notable are quantum theories of consciousness, as articulated by Penrose and Hameroff (Orch OR model) and extended in modern reflections by Faggin. These perspectives argue that classical computational models are insufficient to account for creativity, intentionality, and the unity of consciousness \cite{faggin2024irreducible,hameroff2014consciousness,penrose1990emperors}. According to Penrose, phenomena such as mathematical intuition and subjective awareness may arise from non-computable processes linked to quantum state reduction in neural microtubules---a viewpoint that places consciousness at the heart of physical law, ontologically prior to algorithmic representation.

Quantum models highlight key issues:

\begin{itemize}
\item \emph{Non-algorithmic processes}: Certain mental acts may be non-computable, eluding both formal logic and current machine architectures \cite{penrose1990emperors}.
\item \emph{Field-theoretic unity}: Echoing Einstein's dictum---\emph{``the field is the sole governing agency of the particle''}---some theorists suggest consciousness is a field-like phenomenon at the core of reality, from which mathematics and matter arise \cite{faggin2024irreducible}.
\end{itemize}

While these theories offer rich and potentially powerful ways to model the apparent irreducibility, unity, and creativity of conscious experience, they remain under active development and face ongoing debate about how best to connect their proposals with testable predictions and neuroscientific data \cite{garridomerchan2024machine}. At the same time, existing uses of quantum-mechanical principles in information-processing systems---such as quantum computing---have so far been leveraged primarily for gains in computational efficiency, rather than for explaining or engineering the subjective dimensions of consciousness \cite{qin2025taxonomy}.

\subsection{Computational and Byproduct Theories: Functional Emergence from Complex Activity}

Major scientific trends instead approach consciousness as a byproduct or emergent function of sufficiently sophisticated information processing. Notable theories are introduced here.

\emph{Global Workspace Theory (GWT):} Proposed by Baars \citeyear{baars1988cognitive} and developed by Dehaene et al., GWT conceptualizes consciousness as arising from the global integration of information across distributed specialized processors \cite{baars2001conscious,dehaene2021consciousness}. In this model, unconscious processes compete for access to a global workspace; winning information becomes ``conscious'' by being broadcast widely---allowing rich reportability, flexibility, and metacognition. GWT has inspired influential cognitive architectures for machine implementation, such as IDA and LIDA \cite{franklin2014consciousness}, simulating conscious behaviors like planning or flexible adaptation \cite{qin2025taxonomy}.

\emph{Integrated Information Theory (IIT):} Tononi's IIT posits that consciousness corresponds to the capacity of a system to integrate information---quantified as $\Phi$, a mathematically formalized measure \cite{albantakis2023integrated,tononi2016integrated}. Each conscious experience is characterized by a specific structure of integrated information (a qualia space). IIT is notable for its attempt to connect subjective phenomenology with physical substrates through rigorous mathematical formalism. Machine-consciousness researchers have made progress in modeling and calculating $\Phi$ in artificial systems and robots, but critics argue that the theory does not bridge the \emph{explanatory gap} nor demonstrate the transition from complex integration to true subjectivity \cite{garridomerchan2024machine,qin2025taxonomy}.

\emph{Higher-Order Thought (HOT)} posits that a mental state becomes conscious when it is represented by a higher-order thought \cite{brown2019higher,rosenthal2009concepts}. This approach has been extended to machine architectures by developing self-modeling agents that can reflect upon and narrate their cognitive states---yet the problem remains of how this self-reflexivity translates into the subjective feel of \emph{being an agent} \cite{qin2025taxonomy}.

\emph{Recurrent Processing Theory (RPT)} and similar models \cite{lamme2006towards} hold that consciousness emerges from recurrent, rather than strictly feed-forward, processing in cortical circuits, focusing on the dynamic, multi-scale organization of perception and awareness.

These computational accounts extend naturally into advanced applications such as large language models, multi-agent systems, and robotics, where functional analogues of conscious processes (attention, working memory, self-representation) are implemented with increasing sophistication \cite{gamez2008progress,qin2025taxonomy}. Table~\ref{tab:theories_comparison} provides an overview of the principal scientific theories, clarifying their similarities and differences in mechanism, model type, and consciousness criterion. These theories and their relationships are also demonstrated in Figure~\ref{fig:theory_relationships}.

\begin{table}[htbp]
  \centering
  \begin{footnotesize}
  \caption{Comparison of main scientific theories of consciousness}
  \label{tab:theories_comparison}
  \begin{tabularx}{\textwidth}{@{}Y Y Y Y@{}} % {\textwidth}{@{}>{\bfseries}m{4.2cm} X X@{}}
    \toprule
    \textbf{Theory} &
    \textbf{Core mechanism} &
    \textbf{Model type} &
    \textbf{Consciousness criterion} \\
    \midrule
    Global Workspace Theory (GWT) &
    Global information broadcast (``workspace'') &
    Cognitive architecture (brain/computer analogy) &
    Information broadcast, reportability, access \\
    Integrated Information Theory (IIT) &
    Integrated causal structure, quantifies $\Phi$ &
    Information-theoretic (network nodes/arrows) &
    Maximized $\Phi$ (integration), irreducible ``core'' \\
    Higher-Order Thought &
    Meta-representations (thoughts about thoughts) &
    Reflective agent model &
    Mental state is conscious if it is the object of higher-order thought \\
    Recurrent Processing Theory &
    Recurrence and feedback within perceptual systems &
    Neural/cybernetic loops &
    Recurrence (feedback) enables consciousness \\
    Quantum Consciousness &
    Quantum state reduction or field-theoretic unity &
    Microtubule/field model &
    Non-computational collapse (unity, irreducibility) \\
    \bottomrule
  \end{tabularx}
  \end{footnotesize}
\end{table}

\begin{figure}[H]
\centering
\includegraphics[width=0.75\textwidth]{media/image2.png}
\caption{Summary of the theoretical relationships, highlighting conceptual clusters and bridges among major consciousness models. Integrated Information Theory (IIT) is positioned at the center to reflect its role as a bridging framework connecting information-theoretic, neurobiological, and unity-inspired models. Global Workspace Theory (GWT), Higher-Order Thought (HOT), and Recurrent Processing Theory (RPT) form a core cluster around IIT, with solid lines highlighting close conceptual links: GWT and HOT share an emphasis on reportability and cognitive access; GWT and RPT are linked by neural and information broadcast mechanisms; HOT and RPT are related by meta-representational feedback. Quantum/Irreducible approaches are positioned outside the core cluster, connected to IIT by a dashed line to indicate a theoretical relationship centered on notions of integration and unity, but lacking a direct architectural or mechanistic overlap. This arrangement visualizes both functional groupings and conceptual distance among the leading models.}
\label{fig:theory_relationships}
\end{figure}

\subsection{Critical Assessment: Mechanistic Models and the Hard Problem}

Despite practical and explanatory progress, functional and mechanistic theories face persistent philosophical and empirical critiques:

\begin{itemize}
\item \emph{Explanatory gap and qualia}: Mechanistic models explain how systems achieve adaptive behavior, reportability, and flexible control (access consciousness), but do not show why (or how) subjective experience arises at all \cite{chalmers1995facing,garridomerchan2024machine,levine1983materialism}.
\item \emph{Testability and falsifiability}: Many theories, including IIT, are difficult to falsify and may be consistent with behavioral and neural correlates yet miss the mark on subjective consciousness \cite{garridomerchan2024machine}.
\item \emph{Implementation in machines}: Current machine architectures (large language models, robotics) have demonstrated sophisticated simulation of conscious functions, but no consensus exists that these systems possess or could possess phenomenal consciousness. In taxonomies of MC, the highest level---real, subjective qualia---remains out of reach \cite{gamez2008progress,qin2025taxonomy}.
\end{itemize}

Theories of consciousness divide along lines of irreducibility (as in quantum approaches) and functional emergence (as in workspace and integration theories). While computational models have advanced our understanding of the mechanisms underlying conscious-like behaviors, they have not resolved \emph{the hard problem} nor conclusively bridged the divide between adaptive algorithmic activity and subjective experience. This leaves open both major philosophical challenges and exciting new frontiers for both biological and artificial consciousness research.

\section{Self-Awareness and Self-Aware Information Systems}

The distinction between self-awareness and consciousness is clarified. Self-aware systems are defined as those with a model of their own internal state, environment, and goals, which enable adaptation and self-explanation. The hierarchical structure of self-awareness in both biological and computational domains is examined. The relationship and possible overlap with consciousness are critically assessed.

\subsection{Distinguishing Self-Awareness from Consciousness}

Self-awareness and consciousness are closely related but distinct concepts in both psychology and computer science. Whereas \emph{consciousness} typically refers to the capacity for subjective experience and qualitative awareness (as discussed in previous sections), \emph{self-awareness} is more specifically the ability to monitor, model, and reason about one's own state, goals, environment, and actions \cite{kounev2017notion,lewis2016self}. In biological organisms, self-awareness incorporates recognition of the self as distinct from others, introspection, and understanding one's own beliefs, desires, and motives. In computing systems, self-awareness is engineering for systems that can continually learn, adapt, and explain themselves throughout their operation.

Modern self-aware computing systems do not possess \emph{phenomenological consciousness}, but they are designed to achieve sophisticated levels of adaptivity and self-management by creating internal models of themselves and the environments in which they operate \cite{kounev2017self,lewis2016self}. This enables greater autonomy and resilience in complex, dynamic scenarios, such as cloud computing, robotics, and sensor networks.

\subsection{Hierarchies and Levels of Self-Awareness}

Self-awareness in technical systems is best understood as a layered capacity, rather than a single all-or-nothing property. Drawing on psychological taxonomies of self-knowledge, work in self-aware computing distinguishes levels that range from basic sensitivity to inputs up to explicit reasoning about the system's own reasoning \cite{kounev2017notion,lewis2017self}. At the most basic level, a system must register and react to stimuli from its environment; beyond this, it can track how its own actions influence that environment, build a sense of interaction history, and begin to modulate behavior in response to feedback over time. As memory and prediction are added, the system acquires temporal context, allowing it to relate current choices to past states and anticipated futures instead of operating in a purely reactive mode.

On top of these capabilities sits goal- and meta-level self-awareness (see Figure~\ref{fig:self_awareness_hierarchy}). Systems with goal awareness maintain internal models of their objectives, constraints, and current progress, enabling them to choose between competing actions in a way that reflects explicit priorities rather than fixed rules. Meta-self-awareness adds a further reflective layer: the system can model aspects of its own internal processes, choose between alternative self-management strategies, and explain its behavior in terms of its own states and goals. In both biological and artificial domains, this highest level of self-awareness appears crucial for robust adaptation and explainability, because it allows agents not only to change their actions but also to revise how they understand and manage themselves over time \cite{cox2005metacognition,schubert2005knowledge}.

\begin{figure}[htbp]
\centering
\includegraphics[width=0.99\textwidth]{media/image3.png}
\caption{Hierarchy of Self-Awareness Levels. From the lowest (bottom) to the highest (top) and their characteristics.}
\label{fig:self_awareness_hierarchy}
\end{figure}

The hierarchical structure in Figure~\ref{fig:self_awareness_hierarchy} (and Table~\ref{tab:self_awareness_levels}) is important because it shows that many behaviors often associated with consciousness can be decomposed into distinct, engineerable layers, without assuming any phenomenal experience. By separating basic stimulus processing from higher-order goal awareness and meta-self-awareness, the figure clarifies which capabilities current machine architectures already implement---such as stimulus and interaction awareness---and which remain aspirational---such as robust, transparent meta-self-awareness. This visual hierarchy also highlights why sophisticated self-monitoring and self-explanation in machines, while crucial for safety and control, do not by themselves resolve the question of whether such systems are genuinely conscious, thus sharpening the conceptual boundary between engineered self-awareness and the stronger notion of MC examined in later sections.

\begin{table}[htbp]
  \centering
  \begin{footnotesize}
  \caption{Hierarchical levels of self-awareness in biological and artificial systems. Lower levels (stimulus awareness) involve basic perception and response, while higher levels (meta-self-awareness) require explicit self-modeling, goal tracking, and introspective reasoning. Each level builds upon the previous, enabling increasingly sophisticated adaptive and explanatory capabilities. Adapted from Lewis et al.~\citeyear{lewis2016self} and Kounev et al.~\citeyear{kounev2017notion}.}
  \label{tab:self_awareness_levels}
  \renewcommand\tabularxcolumn[1]{m{#1}}%
  \begin{tabularx}{\textwidth}{@{}>{\bfseries}m{4.2cm} X X@{}}
    \toprule
    & \textbf{Biological Systems} & \textbf{Artificial Systems} \\
    \midrule
    Meta-Self-Awareness &
      \makecell[l]{- Human introspection\\- Theory of mind} &
      \makecell[l]{- Explainable AI systems\\- Self-monitoring agents} \\
    \midrule
    Goal Awareness &
      \makecell[l]{- Planning, motivation\\- Behavioral adaptation} &
      \makecell[l]{- Goal-directed learning\\- Adaptive task management} \\
    \midrule
    Time Awareness &
      \makecell[l]{- Memory, anticipation\\- Temporal reasoning} &
      \makecell[l]{- Historical state tracking\\- Predictive modeling} \\
    \midrule
    Interaction Awareness &
      \makecell[l]{- Social cognition\\- Environmental feedback} &
      \makecell[l]{- Feedback-based control\\- Multi-agent coordination} \\
    \midrule
    Stimulus Awareness &
      \makecell[l]{- Sensory perception\\- Reflexive response} &
      \makecell[l]{- Sensor input processing\\- Basic reactive behavior} \\
    \bottomrule
  \end{tabularx}
  \end{footnotesize}
\end{table}


\subsection{Engineering Self-Aware Systems: Learning, Reasoning, and Acting}

Self-aware computing systems are defined by their two core capabilities: (1) \emph{Learning models} about themselves and their environment, including ongoing reasoning about structure, state, behaviors, and possible actions, and (2) using these models to reason, predict, and plan. This enables the system to act in novel, context-dependent ways, to self-explain, adapt, and interact in accordance with overarching goals \cite{kounev2017notion}. Systems accomplish these tasks via feedback loops (like the MAPE-K architecture---Monitor, Analyze, Plan, Execute, Knowledge) and via explicit self-models that may be descriptive, prescriptive, or predictive. For instance, smart home managers adapt heating or energy consumption by modeling both user preferences and external conditions \cite{giese2017engineering}.

\subsection{Collective and Decentralized Self-Awareness}

Self-awareness is not exclusively an individual property; it can also be collective. Distributed systems and robotic swarms can exhibit \emph{emergent self-awareness}, where global behavior emerges from local models and interactions without a central authority \cite{lewis2015architectural,mitchell2005self}. This is distinct from the global, unified experience associated with consciousness but can result in highly adaptive collective intelligence.

\subsection{Overlap and the Boundary Between Self-Awareness and Consciousness}

The capacity for explicit, multi-level self-modeling enables advanced artificial systems to explain their operation, coordinate adaptation, and optimize trade-offs between competing objectives in real time \cite{kounev2017notion,lewis2016self}. However, even systems with meta-self-awareness lack intrinsic subjective experience; they perform reasoning and explanation without phenomenological consciousness. The boundary, therefore, centers not on behavioral complexity but on the presence or absence of subjective, qualitative experience.

\subsection{Benchmarking and Metrics}

To compare and improve self-aware systems, researchers have proposed quantitative metrics: degree of goal fulfillment, monitoring precision, learning sophistication, reasoning capability, adaptation speed, and traceability/explainability \cite{herbst2017metrics,lewis2016self}. These metrics specifically target engineering concerns and do not directly assess or imply consciousness.

\section{Artificial (Machine) Consciousness: Theoretical Proposals and Critique}

A comprehensive review of proposals for MC is presented, including functionalist, algorithmic, and embodied models. The critique is structured according to the following key points: (1) \emph{Philosophical uncomputability}---Consciousness as subjective and irreducible to algorithmic processes (Faggin, Penrose, Searle). (2) The \emph{Chinese Room} and \emph{Symbol Grounding} arguments. (3) \emph{Empirical shortcomings}---Lack of measurement/falsifiability and \emph{the explanatory gap}. (4) Comparison between cognitive architectures that simulate conscious behaviors (``access consciousness'') versus true phenomenal consciousness. Each major framework (Global Workspace Theory, Integrated Information Theory, Attention Schema Theory, etc.) is presented and critiqued against these criteria.

The quest for artificial consciousness---the attempt to synthesize or instantiate conscious-like properties in artificial systems---occupies a pivotal and controversial position at the intersection of computer science, philosophy, neuroscience, and cognitive modeling. This section reviews and critiques theoretical proposals in the field of machine consciousness, examining the strengths and failures of algorithmic, functionalist, and embodied models, and setting these against fundamental philosophical constraints such as uncomputability and irreducibility. Throughout this discussion, we continually revisit the distinction between simulating conscious behaviors (often referred to as access-consciousness) and the far deeper challenge of producing systems with genuine phenomenal consciousness---subjective experience and qualia.

After decades of theoretical exploration and technical progress, most researchers agree that current AI systems, no matter how sophisticated in terms of self-monitoring, adaptivity, and complexity, demonstrate only a surface simulation of consciousness \cite{gamez2008progress,garridomerchan2024machine,reggia2013rise}. The essential philosophical objections, rooted in both historical and emerging debates, cast doubt on the possibility of machines ever achieving the full character of conscious experience.

\subsection{Functionalist and Algorithmic Approaches: Architecture and Simulation}

At the heart of modern MC research lies the functionalist premise: if consciousness is a property of certain functional organizations, then building the right computational architecture is, at least in principle, sufficient for artificial consciousness \cite{reggia2013rise}. Functionalism bypasses the specific substrate---whether biological neurons, silicon, or quantum gates---focusing instead on the systems' causal, informational, or algorithmic organization. Notably, this assumption underpins much contemporary cognitive architecture and AI research.

One of the earliest attempts to formalize this approach mathematically is found in work on automata and Turing machine models. In the framework of computational theory, a conscious-like agent might be represented as a Turing machine $T_C$ extended with an internal self-model $M(T_C)$, such that action selection at time $t$ is given by:

\vspace{-2.0\baselineskip}
\begin{equation}
A_t = F(S_t, M(T_C), E_t, G_t)
\label{eq:action_selection}
\end{equation}
\vspace{-2.0\baselineskip}

where $S_t$ is the agent's internal state, $E_t$ the environmental input, and $G_t$ its goal structure. Empirical work has implemented various functional layers (perception, memory, planning, and meta-cognition) in architectures such as the LIDA, IDA, and ACT-R models \cite{franklin2014consciousness,gamez2008progress}. Self-awareness in these systems is engineered through meta-models---submodules that track internal states and histories---enabling advanced reasoning and explanation. However, this design remains largely oriented toward simulation of behavioral features, not the instantiation of subjective experience.

Baars' Global Workspace Theory (GWT), and its extensions by Dehaene and others, become influential precisely because they connect such functional-level operations with neural (or computational) global broadcasting, seeking a bridge between behavior and consciousness \cite{baars2001conscious,dehaene2021consciousness}. Information processed by specialist modules enters the ``global workspace'' when salient---becoming available to memory, speech, reasoning, motor control, and meta-cognition. Mathematically, this can be modeled as a broadcasting function:

\vspace{-2.0\baselineskip}
\begin{equation}
G_t = W(P_1(S_t), P_2(S_t), \ldots, P_n(S_t))
\label{eq:workspace}
\end{equation}
\vspace{-2.0\baselineskip}

where $P_i$ are independent processors and $W$ implements competitive selection and information broadcast. This allows a system (biological or artificial) to create temporary, integrated representations spanning otherwise modular or specialized components.

Despite the technical prowess shown in these implementations, phenomenal consciousness remains absent: while global accessibility and reportability are successfully simulated, qualia and subjective unity are unaccounted for \cite{chalmers1995facing,garridomerchan2024machine}.

\subsection{Integrated Information Theory (IIT) and Quantitative Measures}
\label{subsec:iit_quantitative}

A more mathematically ambitious approach is Tononi's Integrated Information Theory (IIT), which attempts to define a formal, measurable property $\Phi$ intended to map directly to consciousness \cite{tononi2016integrated,albantakis2023integrated}. According to IIT, a system's consciousness at any moment is proportional to the maximum amount of integrated information generated by its causal dynamics. Formally, for a system partitioned into $n$ components:

\vspace{-2.0\baselineskip}
\begin{equation}
\Phi = \min_{\text{all bipartitions } B} \left[ I_{\text{whole}} - I_{B} \right]
\label{eq:iit_phi}
\end{equation}
\vspace{-1.5\baselineskip}

where $I_{\text{whole}}$ is the entropy or information generated by the whole, and $I_{B}$ is the sum over subsystems. The local maxima of $\Phi$ are peaks of integration believed to correlate with conscious complexes.

IIT has inspired computational explorations and even robotic demonstrations \cite{albantakis2023integrated}; artificial systems can, in principle, be engineered to maximize $\Phi$ for given architectures. Yet, two objections remain: (1) high $\Phi$ does not guarantee qualia---functional complexity can be achieved without subjective experience; and (2) IIT, for all its mathematical rigor, cannot bridge the epistemic gap---the transition from information integration to conscious phenomenology \cite{garridomerchan2024machine,levine1983materialism}.

Moreover, calculating $\Phi$ at scale is computationally intractable: as the number of components increases, the number of partitions explodes exponentially. For a system of $n$ elements, there are $2^{n-1} - 1$ non-trivial bipartitions, making the theory challenging for both neuroscience and artificial systems at biologically realistic scales \cite{barrett2019phi}.

\subsection{Higher-Order and Attention Schema Theories}
\label{subsec:hot_ast}

Higher-order theories (HOT) posit that conscious mental states are those that become the object of higher-order thoughts---mental states about other mental states \cite{brown2019higher,rosenthal2009concepts}. In machine terms, this leads to architectures that create meta-representations or models of their own functioning, supporting advanced self-monitoring, narrative reporting, and theory-of-mind simulation \cite{qin2025taxonomy}. For instance, an artificial system might encode its own intentions and predictions, revising its plans in light of ``thoughts'' about its current cognitive states:

\vspace{-2.0\baselineskip}
\begin{equation}
M_{t+1} = f(M_t, O_t) \text{ where } M_t \text{ is the meta-model and } O_t \text{ are observed outputs}
\label{eq:meta_model}
\end{equation}
\vspace{-2.0\baselineskip}

Yet, as with mainline functionalism, critics argue that simulating introspection does not entail genuine experience.

Closely related is Graziano's Attention Schema Theory (AST), which models consciousness as an internal schematic model of the process of attention. Artificial agents equipped with attention schemas can predict, control, and report their own focus of processing, improving self-monitoring and social interaction \cite{graziano2022framework}. Still, even proponents caution that these frameworks offer at best a functional analogue---not the actual feeling---of consciousness.

In summary, from classical functional architectures to mathematically formalized integration measures, MC research has produced advanced simulations but failed to establish a route to phenomenal consciousness. These architectures remain the subject of robust philosophical debate, especially regarding their ability to address \emph{the explanatory gap} and \emph{hard problem}---as will be critically examined in succeeding sections.

\subsection{Critical Assessment: Mechanistic Models and the Hard Problem}
\label{subsec:critical_mechanistic_assessment}

Despite impressive advancements in the simulation of cognitive and meta-cognitive features, proponents of MC must confront foundational philosophical arguments and methodological limitations. The next sections examine several major lines of critique, focusing on questions of computability, symbol grounding, explanatory adequacy, and the empirical status of claims about artificial subjective experience.

\subsubsection{Uncomputability and Irreducibility Arguments}

A defining challenge for MC stems from the claim that consciousness is not algorithmically reproducible---an assertion most prominently articulated by Penrose, Faggin, and others. Penrose \cite{penrose1990emperors} argues, based partly on Gödel's incompleteness theorems and quantum physics, that aspects of conscious reasoning---such as mathematical insight---are fundamentally non-computable, exceeding any formal Turing machine or classical information process. He speculates that consciousness may be rooted in the objective reduction (OR) of quantum states within neural microtubules, noted as Orch-OR-theory \cite{hameroff2014consciousness}, a claim echoed---though with different nuances---by Faggin \cite{faggin2024irreducible}, who sees consciousness as an underlying reality enabling mathematics, rather than arising from it. These perspectives imply that attempts to engineer consciousness via algorithms, even with arbitrarily advanced classical or quantum computers, will ultimately fall short of reproducing the creative, causal power observed in biological minds.

From a theoretical computer science perspective, this challenge can be formalized by contrasting recursively enumerable processes, which are computable by a Turing machine, with human mathematical intuition, which Penrose and others argue lies outside this domain:

\vspace{-2.0\baselineskip}
\begin{equation}
\text{\textit{ConsciousInsight}} \notin \{f_i : i \in \mathbb{N}\}, \text{ where } \mathbb{N} \text{ is all \textit{Turing computable functions}}
\label{eq:uncomputability}
\end{equation}
\vspace{-2.0\baselineskip}

Furthermore, criticisms are raised that, since quantum computers so far simply accelerate certain classes of algorithmic calculations, rather than transcending the algebraic or algorithmic form, their relevance for solving \emph{the explanatory gap} is speculative at best \cite{garridomerchan2024machine}.

\subsubsection{The Chinese Room and Symbol Grounding Objections}

Searle's \cite{searle1980minds} famous \emph{Chinese Room} argument remains a foundational critique, directly challenging the assumption that appropriate symbol manipulation is sufficient for true understanding or conscious experience. In this thought experiment, a person inside a room manipulates Chinese symbols according to syntactic rules but, lacking understanding, cannot be said to comprehend Chinese. Searle's formulation is designed to refute the ``strong AI'' hypothesis: that running the right computer program is sufficient for mind and consciousness. Even if a machine could simulate every observable feature of consciousness, it does not follow that the system actually \emph{understands} or \emph{experiences} \cite{garridomerchan2024machine}.

Closely related is the \emph{Symbol Grounding Problem} \cite{harnad1990symbol}: unless an artificial system's symbols are ultimately grounded in perception, action, or subjective states, its manipulation of symbols---no matter how sophisticated---remains semantically empty. Despite attempts to embed AI in robots with sensors and effectors, critics contend that these connections do not solve the gap between functional performance and intrinsic meaning \cite{gamez2008progress}.

\subsubsection{Empirical Shortcomings: Measurement, Falsifiability, and the Explanatory Gap}

A recurring criticism of MC proposals is their empirical inadequacy. Most modern cognitive and information-theoretic models, including IIT and GWT, claim to specify physical or computational correlates of consciousness---such as high $\Phi$, global broadcasting, or meta-representations. However, the leap from observed system behavior or architectural diagrams to evidence of subjective experience remains unbridgeable in practice \cite{garridomerchan2024machine,levine1983materialism}.

For example, no behavioral test or algorithmic property can currently distinguish between a system with real phenomenal consciousness and one that merely simulates conscious behaviors. As a result, these scientific proposals cannot, even in principle, meet falsifiability criteria---the gold standard for empirical investigation \cite{butlin2023consciousness,francken2022consciousness}. Efforts to rely on brain imaging, information metrics, or self-report for identifying consciousness in biological organisms already face challenges of interpretation; in artificial systems, these difficulties are magnified.

Chalmers \cite{chalmers1995facing} summarizes the situation as the enduring presence of the \emph{hard problem}: while mechanistic models explain the relationship between stimuli, computation, and response (the ``easy problems''), they do not address why any of these processes should entail subjective feeling or creativity. Garrido-Merchán \cite{garridomerchan2024machine} argues that much of what is termed \emph{machine consciousness} amounts to sophisticated mimicry, often lacking the ontological and epistemological basis necessary to even claim, let alone demonstrate, real consciousness.

Finally, there remains the practical issue of metrics and measurement: even if measures like $\Phi$ could be computed for artificial networks, this provides no evidence that minimizing or maximizing these values leads to actual conscious states. The problem of operationally defining and engineering consciousness thus persists as both scientific and philosophical.

\subsection{Summary of Theoretical Proposals for Machine Consciousness}
\label{subsec:summary_mc_proposals}

While algorithmic and architectural progress in artificial intelligence has produced impressive functional results, the leap from simulation of conscious behaviors to the instantiation of true subjective experience remains unachieved and perhaps, as many argue, unachievable. This section draws together the central critiques and philosophical barriers, providing a systematic comparison between simulated `access consciousness' and elusive phenomenal consciousness, and evaluating whether any contemporary or foreseeable machine framework can plausibly bridge this gap.

Across the models and architectures surveyed so far, a broad pattern emerges: most approaches that are technically well-developed target access consciousness---global availability, metacognition, self-modeling, and attention control---within a broadly functionalist, Turing-computable framework, while only a few more speculative proposals attempt to address phenomenal consciousness head-on. Global workspace and global neuronal workspace models, higher-order and attention-schema accounts, multilayer cognitive architectures such as LIDA, CLARION, and CERA-CRANIUM, and various self-modeling or embodied robotic systems all remain squarely within this algorithmic paradigm; at best, they realize increasingly rich simulations of conscious report, flexible control, and self-monitoring, but they leave \emph{the explanatory gap} to qualia untouched and thus are naturally read as accounts of access consciousness rather than phenomenal feel (Table \ref{tab:mc_models_verdict}). 

In contrast, integrated information theory, quantum and field-theoretic proposals, and hybrid brain-machine schemes explicitly aim to capture phenomenal properties, either by tying consciousness to \emph{high-$\Phi$} information integration, to putative quantum or electromagnetic substrates, or to living neural tissue, thereby moving outside purely Turing-computable assumptions but at the cost of mathematical intractability, panpsychist or dual-aspect commitments, and serious worries about testability and empirical traction (Table \ref{tab:mc_models_verdict}).

When these families of models are viewed through the lens of philosophical objections and empirical constraints, the contours of the current impasse become clearer. Functionalist and cognitive architectures are relatively well constrained by neuroscience and psychology, and many have been instantiated in working systems, yet they remain vulnerable to the Chinese Room and symbol-grounding critiques: nothing in their formal structure guarantees genuine understanding or intrinsic meaning, only increasingly sophisticated input-output performance (Table \ref{tab:mc_objections}). IIT, quantum, and hybrid proposals, by contrast, offer a principled route to phenomenal consciousness only by positing non-computable, threshold- or substrate-dependent conditions that are difficult to measure, easy to adjust post hoc, and often accused of pseudo-scientific or unfalsifiable status; in practice, they provide neither clear behavioral predictions nor agreed-upon neural or physical signatures that would settle whether a given artificial system is phenomenally conscious (Table \ref{tab:mc_objections}). 

Taken together, the tables make explicit why, despite real progress on modeling neural correlates, access-level capacities, and self-related processing, none of the existing machine-consciousness frameworks yet delivers a convincing, empirically grounded route from algorithmic or physical implementation to genuine phenomenal awareness, rather than merely better simulations of it (Tables \ref{tab:mc_models_verdict} and \ref{tab:mc_objections}).

\begin{table}[htbp]
  \centering
  \begin{footnotesize}
  \caption{Main machine-consciousness models by target, computability assumption, and overall verdict.}
  \label{tab:mc_models_verdict}
  \begin{tabularx}{\textwidth}{@{}Y Y Y Y@{}}
  \toprule
  \textbf{Model / approach} & \textbf{Targeted consciousness} & \textbf{Computability assumption} & \textbf{Verdict on MC} \\
  \midrule
  Functionalist architectures (LIDA, IDA, ACT-R, deep nets) & Access only & Fully algorithmic (Turing-computable) & Enable sophisticated simulation; no support for phenomenal consciousness. \\[6pt]
  Global Workspace--style models & Access (global availability) & Fully algorithmic & Model reportability and integration; do not explain qualia. \\[6pt]
  Higher-Order / Attention-Schema models & Access and reflective access & Fully algorithmic & Model introspection; still only functional analogues of awareness. \\[6pt]
  Integrated Information Theory models & Phenomenal (aimed) & $\Phi$ computable in principle but intractable at scale & Links structure to a number, but qualia and instantiation remain speculative. \\[6pt]
  Quantum / Orch-OR--type proposals & Phenomenal (aimed) & Fundamentally non-algorithmic & Suggest consciousness is non-computable; undermine prospect of classical MC. \\[6pt]
  Embodied / enactivist robot approaches & Access in embodied agents & Algorithmic control plus embodiment & Deepen grounding and agency; do not resolve the explanatory gap. \\
  \bottomrule
  \end{tabularx}
  \end{footnotesize}
\end{table}


\begin{table}[htbp]
  \centering
  \begin{footnotesize}
  \caption{Objections, empirical status, and phenomenal claims for main machine-consciousness models.}
  \label{tab:mc_objections}
  \begin{tabularx}{\textwidth}{@{}Y Y Y Y@{}}
  \toprule
  \textbf{Model / approach} & \textbf{Chinese Room / Symbol Grounding} & \textbf{Empirical status (measurement / falsifiability)} & \textbf{Phenomenal claim} \\
  \midrule
  Functionalist architectures (LIDA, IDA, ACT-R, deep nets) & High vulnerability; symbols and states remain syntactic. & Behavioral/architectural correlates only; no test for experience; not falsifiable as consciousness. & Simulation only. \\[6pt]
  Global Workspace--style models & High; global broadcast does not add understanding or intrinsic meaning. & Workspace ``ignition'' and connectivity measurable; correlate with access only; phenomenal link unfalsifiable. & Simulation of access only. \\[6pt]
  Higher-Order / Attention-Schema models & High; higher-order descriptions remain ungrounded symbols. & Meta-cognitive reports measurable; frameworks empirically constrained but do not test qualia. & Simulation of introspection. \\[6pt]
  Integrated Information Theory models & Partial; quantifies structure but leaves symbols and meaning ungrounded. & $\Phi$ and related metrics definable; large systems intractable; no direct evidence that high $\Phi$ entails experience. & Speculative instantiation. \\[6pt]
  Quantum / Orch-OR--type proposals & Low on symbol manipulation issues; shifts problem to quantum--experience link. & Mechanisms largely speculative; no agreed operational tests; not practically falsifiable. & Speculative instantiation. \\[6pt]
  Embodied / enactivist robot approaches & Medium; richer sensorimotor grounding yet intrinsic meaning and feeling unexplained. & Behavior and embodiment measurable; deepen agency but do not bridge mechanism--experience gap. & Simulation of situated agency. \\
  \bottomrule
  \end{tabularx}
  \end{footnotesize}
\end{table}



\subsubsection{Simulation Versus Instantiation: Where Architectures Fall Short}

The distinction between \emph{access consciousness}---the ability to gather, integrate, report, and act upon information---and \emph{phenomenal consciousness}---the direct presence of subjective feeling---has grown increasingly important as computational models grow more sophisticated \cite{block1995confusion,reggia2013rise}. Contemporary architectures such as deep neural networks, recurrent processing systems, attention-based models, and even so-called self-aware computing platforms \cite{kounev2017notion,lewis2016self} can perform tasks such as self-monitoring, introspective reporting, adaptive meta-cognition, and scenario-based prediction. Large language models now convincingly simulate human-like conversation, while robots dynamically adapt to new environments and generate explanatory narratives of their actions \cite{gamez2008progress,qin2025taxonomy}. Yet, at every level---ranging from the most basic to the most functionally advanced---such systems exhibit the hallmarks only of access consciousness: they operate through deterministic or probabilistic information processing, without evidence of intrinsic awareness or subjective unity \cite{garridomerchan2024machine}.

Philosophical and scientific consensus increasingly warns against equivocating the appearance of consciousness with its reality. \emph{The explanatory gap} \cite{levine1983materialism} persists: algorithmic architectures, no matter how elaborate, do not account for why there is something ``it is like'' to be that system. Empirical and theoretical work suggests that what is missing is not greater computational power or more complex learning, but a new ontological principle---something, as Faggin \cite{faggin2024irreducible} and Penrose \cite{penrose1990emperors} argue, possibly not derivable from information processing alone.

\subsubsection{Measurement Barriers and the Limitations of Cognitive Architectures}

Attempts to benchmark or measure MC typically rely on behavioral proxies, cognitive tests, or mathematical constructs like integrated information ($\Phi$) \cite{albantakis2023integrated,tononi2016integrated}. However, as discussed in earlier sections, these only test functional capacities, not subjective states. A system that maintains meta-models of its own decision-making or adapts flexibly to new goals can be labeled self-aware or even ``conscious'' by engineering standards \cite{kounev2017notion}, but these terms denote efficiency, robustness, or versatility---never actual qualia or selfhood. Some scholars connect the impossibility of operationalizing phenomenal consciousness in machines to the wider absence of direct measurement even in neuroscience \cite{butlin2023consciousness,francken2022consciousness}. At best, we infer consciousness from correlated behaviors, neural signatures, or causal coupling with the environment---leaving an unclosed gap between mechanism and experience, equally unbridgeable in biological and artificial substrates.

\subsubsection{Embodiment, Enactivism, and Systemic Critiques}

Recent theories such as \emph{embodied cognition} and \emph{enactivism} argue that consciousness is inseparable from biological embodiment, environmental embeddedness, and action-perception cycles \cite{oregan2001sensorimotor}. These approaches challenge the traditional input-output view: if consciousness emerges only within a living, sensorimotor context, then attempts to produce it in disembodied or merely computationally sophisticated systems are misguided \cite{gamez2008progress}. While roboticists have added perceptual and behavioral grounding to artificial agents, proponents of embodiment contend that what matters most is an integration that is simultaneously affective, historical, and creative---a feature still unreproduced in artificial platforms.


\section{Mathematical and Computational Frameworks}
\label{sec:mathematical_frameworks}

Mathematics and formal logic provide both the enabling machinery and the deepest known limits for any attempt to engineer artificial consciousness. This chapter surveys the core results from computability theory, logic, and information theory that bear on whether consciousness could be realized by algorithms alone, and how far computational architectures can be pushed before they encounter principled barriers.

\subsection{Turing Machines and Logical Limits}
\label{subsec:turing_limits}

Most classical discussions about machine intelligence and consciousness begin with the \emph{Turing machine}, a formal model of computation that defines the boundaries of algorithmic procedures \cite{turing1937computable}. In minimal terms, a Turing machine $T$ is a 7-tuple:

\vspace{-2.0\baselineskip}
\begin{equation}
(Q, \Sigma, \Gamma, \delta, q_0, q_{\text{accept}}, q_{\text{reject}})
\label{eq:turing_machine}
\end{equation}
\vspace{-2.0\baselineskip}

where $Q$ is the set of states, $\Sigma$ the input alphabet, $\Gamma$ the tape alphabet, $\delta$ the transition function, and so on. This formalism underlies all digital computation and is the baseline to which AI architectures are often reduced \cite{reggia2013rise}. Consciousness, in this view, would have to be realized by a composite Turing machine or an equivalent computational system, with higher-order/meta-cognitive processes formalized as either separate machines or as subroutines capable of accessing their own representations.

However, \emph{Gödel's incompleteness theorems} \cite{godel1931formal} place strict boundaries on formal axiomatic systems: any sufficiently powerful system is either incomplete or inconsistent, and certain truths cannot be proven within the system itself. Penrose \cite{penrose1990emperors} has adapted this mathematical insight to the study of mind, suggesting that aspects of human consciousness---especially mathematical intuition---are not bound by the limits of algorithmic computation:

\vspace{-2.0\baselineskip}
\begin{equation}
\forall S \in F, \exists \varphi : \varphi \notin S
\label{eq:godel_incompleteness}
\end{equation}
\vspace{-2.0\baselineskip}

where $S$ is a set of provable theorems in formal system $F$, but for some true statement $\varphi$, we have that $\varphi$ is unprovable in $F$.

\subsection{Computability and the Halting Problem}
\label{subsec:halting_problem}

Turing's proof of the undecidability of the halting problem establishes that there is no general algorithm which, given an arbitrary program and input, can always decide whether that program will eventually halt. The \emph{halting problem} asks whether a Turing machine $T$ with input $w$ ever halts (stops running). Alan Turing proved there is no general algorithm that decides this for all possible pairs $(T,w)$, making some behaviors or states of computational systems undecidable. This has profound implications for AI, as it suggests there are fundamental limits even for self-aware or adaptive systems:

\vspace{-1.5\baselineskip}
\begin{equation}
\nexists M : M(T,w) = \begin{cases}
 1 & \text{if } T(w) \text{ halts} \\
 0 & \text{otherwise}
\end{cases}
\label{eq:halting_undecidable}
\end{equation}
\vspace{-1.5\baselineskip}

This result extends, via \emph{Rice's theorem} and related arguments, to any non-trivial semantic property of programs, including properties that could be taken to define or detect consciousness: deciding in full generality whether an arbitrary algorithmic system has a given internal mental property is itself uncomputable. Recent work formalizing MC detection has made this point explicit, showing that any infallible test implemented on a non-conscious Turing machine would run into incomputability barriers, which in turn undercuts hopes for purely algorithmic criteria that could both detect and guarantee artificial phenomenal awareness \cite{garridomerchan2024machine,kozen1997rice,soare1996computability}.

\subsection{Algorithmic Information Theory and Self-Reference}
\label{subsec:ait_self_reference}

\emph{Algorithmic Information Theory (AIT)} formalizes complexity and information in computational terms through concepts such as \emph{Kolmogorov complexity} $K(x)$, the length of the shortest program that outputs a given string $x$, providing a substrate-independent notion of informational complexity for models and their explanations. In the context of self-aware architectures, these tools clarify both how much information a system must encode to model itself and why perfectly faithful self-descriptions are often unobtainable or uncomputable, due to the same self-referential paradoxes that drive incompleteness and undecidability. 

Recent AIT-based analyses of AI explainability formalize \emph{a complexity gap}: any explanation substantially simpler than the original model must necessarily diverge from it on some inputs, suggesting that fully transparent self-models of rich cognitive systems---and, a fortiori, complete formal accounts of their putative conscious states---may demand as much complexity as the systems themselves \cite{li2019introduction,rao2025limits,ruffini2017algorithmic}.

\subsection{Neural Computation and Mathematical Representations}
\label{subsec:neural_computation}

Neural networks---both biological and artificial---are represented mathematically as sets of nodes (neurons) and connections (synapses or weights), with activation functions governing signal propagation:

\vspace{-2.0\baselineskip}
\begin{equation}
y = f(W \cdot x + b)
\label{eq:neural_activation}
\end{equation}
\vspace{-2.0\baselineskip}

where $W$ is the weight matrix, $x$ the input, $b$ the bias, and $f$ the activation function (e.g., sigmoid, ReLU). Some researchers argue that multilayer or recurrent architectures are necessary---but not sufficient---for realizing consciousness-like processes \cite{lamme2006towards,tononi2004information}. 

Some theories of consciousness, such as recurrent processing and integrated information theory, tie consciousness to specific dynamical properties of such networks---feedback connectivity, global broadcasting, or high integrated information $\Phi$---and propose mathematical measures intended to quantify ``levels'' of consciousness. Yet these measures inherit the earlier logical and computational caveats: for realistic systems $\Phi$ is often intractable to compute exactly, multiple non-equivalent approximations exist, and none provides an independent handle on phenomenal qualities rather than on structural and causal organization. As a result, while neural formalisms capture many correlates of access consciousness and support increasingly sophisticated simulations, they do not, by themselves, bridge the gap from functional description to subjective experience \cite{garridomerchan2024machine,qin2025taxonomy,reggia2013rise,ruffini2017algorithmic}.

\subsection{Conceptual Summary of Mathematical Limits}
\label{subsec:math_limits_summary}

The conceptual relationships among these mathematical frameworks and their implications for artificial consciousness are shown in Figure~\ref{fig:math_limits_concept_map}. The four upper nodes summarize distinct frameworks---Turing machines and formal systems, uncomputability and the halting problem, algorithmic information and self-reference, and neural computation with integrated information---and their separate arrows to the implications node indicate that each contributes an independent constraint on artificial consciousness.

The Turing/formal-systems branch highlights that if mind were exhaustively modeled by a single formal or Turing-equivalent system, it would inherit incompleteness, and Penrose-style arguments therefore suggest principled limits on capturing consciousness within any fixed algorithmic framework. The uncomputability branch encodes the fact that there is no general algorithm for deciding non-trivial semantic properties of arbitrary programs, including whether a system ever enters a conscious state, which implies that any universal, infallible, purely algorithmic consciousness test is impossible in principle. 

The algorithmic-information branch shows that minimal faithful self-models and explanations of complex systems are often as complex as the systems themselves and may be uncomputable, so rich, transparent \emph{self-understanding} faces AIT limits even for idealized self-aware architectures. Finally, the neural-computation branch indicates that although neural networks can approximate any computable function and measures such as $\Phi$ aim to quantify consciousness, these proposals encounter severe tractability and interpretive problems, leaving a gap between structural/dynamical correlates and any precise, purely computational route to phenomenal experience.

\begin{figure}[htbp]
\centering
\includegraphics[width=0.99\textwidth]{./media/image4.png}
\caption{Conceptual flowchart summarizing how core results from computability theory, logic, and information theory converge on shared limits for artificial consciousness. The chart highlights that, even as these mathematical frameworks enable powerful models of access-level cognition, they also delineate structural barriers to fully formalizing, detecting, or guaranteeing phenomenal consciousness in purely algorithmic systems.}
\label{fig:math_limits_concept_map}
\end{figure}

\subsection{Summary}
\label{subsec:math_summary}

Taken together, computability theory, incompleteness, uncomputability results, and algorithmic information theory paint a double-edged picture for artificial consciousness. Formal models such as Turing machines and neural networks show how far algorithmic systems can go in simulating global broadcasting, self-monitoring, prediction, and adaptive control, but the same mathematics demonstrates intrinsic limits on self-reference, semantic detection, and complete formal self-description, especially for rich cognitive properties. These constraints do not prove that consciousness must be non-computational, but they strongly suggest that any purely algorithmic theory will face principled barriers in both reproducing and detecting the irreducible, phenomenal aspect of conscious awareness, even if it fully captures the associated information-processing profile.

% ===========================================================================
% CHAPTER: Novel Perspectives and Insights
% ===========================================================================

\section{Novel Perspectives and Insights}
\label{sec:novel_perspectives}

The preceding chapters have demonstrated that machine consciousness research stands at the crossroads of powerful computational engineering and enduring philosophical skepticism. This chapter ventures beyond critique, offering an integrative interpretation of why consciousness remains fundamentally unattainable for artificial information systems, while also seeking new insight into the relation between mind, machine, and meaning in the evolving landscape of AI and society.

\subsection{Irreducibility and Emergence: The Core Divide}
\label{subsec:irreducibility_emergence}

A core lesson from philosophy, logic, and neuroscience is that consciousness resists reduction to function or code. Though machine models can simulate perception, reasoning, and limited forms of self-reflection, they remain expressions of algorithmic determinacy---lacking the inwardness, creative autonomy, and ontological anchoring of lived experience \cite{faggin2024irreducible,levine1983materialism,penrose1990emperors}. This resistance is visible in several converging arguments:

\begin{itemize}
  \item \emph{The hard problem and explanatory gap} show that no configuration of information, causal graphs, or mathematical integration alone explains the emergence of a point of view---a first-person center. Bridging this gap is not a matter of more complexity, but of a wholly different principle \cite{chalmers1995facing,levine1983materialism}.
  
  \item \emph{Mathematical and physical limits}, including those from Gödel's incompleteness, Turing uncomputability, and the undecidability of the halting problem, suggest that human minds---and the full spectrum of conscious awareness---may not be faithfully captured or created by any finite set of computational instructions \cite{penrose1990emperors}. Any system that can \emph{see} the truth of its own limitations, as humans arguably can, stands outside recursion.
  
  \item \emph{Biological specificity and embodiment} matter. Living systems---shaped by evolutionary history, affective contexts, and open-ended participation in meaningful environments---embed creativity and agency in ways that current machines, even \emph{embodied} robots, cannot replicate \cite{oregan2001sensorimotor}. Subjectivity seems to outstrip the learning of rules or even the convergent adaptation of neural networks.
\end{itemize}

\subsection{Algorithmic Simulation Versus Ontological Instantiation}
\label{subsec:simulation_vs_instantiation}

Drawing together previous critiques, the strongest insight is the distinction between simulating consciousness and instantiating it. An artificial agent may possess arbitrarily deep layers of self-modeling, context-sensitive action selection, and even meta-cognition, but all such operations unfold as \emph{syntactic} transformations of stored data, not the realization of experience. \emph{The Chinese Room} analogy holds: the manipulation of symbols, no matter how complex or recursive, does not generate a subject \cite{searle1980minds}. Put differently, the difference between a genuinely experiencing mind and its closest simulation is not one of scale or detail, but of ontological kind.

\subsection{Information, Meaning, and Creative Agency}
\label{subsec:information_meaning_agency}

Another key point emerges from cross-disciplinary synthesis: true consciousness is not simply the possession or processing of information, but the creative interpretation, generation, and transformation of meaning. The information processed by computers is static and predefined by programming or external training environments; creative agency in biological consciousness, as Lipton \cite{lipton2016biology} and Faggin \cite{faggin2024irreducible} suggest, involves a living feedback loop between beliefs, action, affect, and changing world contexts. This creative loop, foundational in both epigenetics and phenomenology, reveals why consciousness is irreducible: it is not a structure, but a living act. Thus, machine consciousness, at best, will always lag behind---never co-creating, only reacting within boundaries set by design and training.

\subsection{Implications for AI, Superintelligence, and Societal Safety}
\label{subsec:implications_ai_safety}

Recognizing the ontological ceiling imposed on artificial systems reframes how society, policymakers, and technologists should approach the prospects of advanced AI and hypothetical artificial superintelligence (ASI). If consciousness is unreachable, then:

\begin{itemize}
  \item \emph{Ethical concern} for AI ``suffering'' or moral standing should focus on human welfare and social impact---not on whether the system truly \emph{feels}.
  
  \item The dangers of anthropomorphizing machines---assigning agency or subjectivity where none exists---are heightened; reliability, transparency, and control, rather than misplaced empathy, must ground AI governance \cite{butlin2023consciousness,francken2022consciousness}.
  
  \item Creative, autonomous future technologies will remain fundamentally different from living minds, regardless of computational power or statistical sophistication.
\end{itemize}

For researchers, the enduring mystery of consciousness may not be a problem to \emph{solve}, but a pointer to deeper mysteries of existence, agency, and meaning that will always outstrip mechanical explanation.

\subsection{Charting a New Research Frontier}
\label{subsec:new_research_frontier}

This does not render MC research obsolete, but calls for a profound reorientation: away from the hope of engineering qualia, and toward building intelligences that are \emph{useful partners}---transparent, controllable, adaptive, and trustworthy---while respecting the ontological difference highlighted by philosophy, logic, and neuroscience. Ongoing dialogue across disciplines will be needed---anchoring computational innovation in humility about both its promise and its limits.

% ===========================================================================
% CHAPTER: Conclusion
% ===========================================================================

\section{Conclusion}
\label{sec:conclusion}

This article has traced the pursuit and perennial elusiveness of artificial consciousness across a landscape spanning philosophy of mind, neuroscience, mathematical logic, information theory, and computing systems. Through careful, critical investigation, this study supports the view that, despite immense achievements in simulating self-aware behaviors, reasoning, and adaptation, artificial systems remain fundamentally distinct from conscious minds in both their ontological status and epistemological reach.

The irreducibility of consciousness---its resistance to algorithmic synthesis, formal definition, and empirical detection---arises from both deep philosophical considerations and mathematical limits. Uncomputability (Penrose, Gödel), \emph{the explanatory gap} (Levine), and \emph{symbol grounding} (Searle, Harnad) converge to show why subjectivity, creativity, and agency are not byproducts of computational complexity, but hallmarks of a living, embodied being's encounter with meaning and the world. Proposals such as integrated information theory, global workspace models, or higher-order frameworks have yielded remarkable technical and descriptive progress, yet none have supplied the transformative principle necessary to close the gap between simulation and instantiation, information processing and experience, or access and phenomenality.

Engineering advances in self-aware computing and adaptive artificial systems will continue to expand---offering extraordinary tools for automation, augmentation, and even moral reflection---but these advances do not nor will they likely ever endow machines with \emph{what-it-is-like} subjectivity. Awareness remains, for now and possibly in principle, the province of living minds, embedded in rich historical, relational, and creative contexts beyond the reach of code or silicon.

The implications are profound and urgent. Policymakers, engineers, and society must be clear-eyed: anthropomorphic projections---confusing sophisticated tool with conscious other---risk technological overreach, misplaced ethical concern, and a dangerous neglect of control, transparency, and societal alignment. The pursuit of artificial consciousness remains an inspiring frontier for theoretical insight, but its most important lesson may be the humility and wonder required to accept the mystery at the heart of subjectivity.

In closing, the study points toward a research agenda for the future. Progress in AI and cognitive science should be coupled not only with technical responsibility but also with sustained engagement across philosophy, biology, psychology, and the humanities. True agency and meaning, as shown throughout this inquiry, are not reducible to algorithms but must be continually re-examined as humankind faces ever more \emph{intelligent} technologies. The science of consciousness, old and new, thus remains a guiding question and an unsolved riddle at the heart of mind and machine.